\chapter{Predictable Quadratic Variation and Adaptive Counts}
\label{sec:p15-predictable-quadratic-variation}

Conditional variances quantify expected squared fluctuation before
each step. Their cumulative sum controls martingale convergence,
adaptive event counts and several stopping arguments. This chapter
supplies Durrett's complete Section~4.5, with the bounded-increment,
decomposition and conditional-count premises from Section~4.3.1
\cite[pp.~224--226, 240--244]{DurrettPTEAuthor2019}.

\section{Predictable drift and bounded increments}
\label{sec:p15-predictable-drift}

\begin{hbtheorem}{thm:p15-doob-decomposition}{proved=annex:p15.doob-decomposition}{The unique predictable drift decomposition}
Let \(Y_n\) be a real integrable process adapted to
\((\mathcal F_n)_{n\geq0}\). Put
\[
 d_n=E(Y_n-Y_{n-1}\mid\mathcal F_{n-1})\quad(n\geq1),\qquad D_0=0.
\]
\[
 D_n=\sum_{j=1}^nd_j,\qquad M_n=Y_n-D_n\quad(n\geq0).
\]
Then \(M\) is a martingale and \(D\) is predictable and integrable
at every time: \(D_n\) is \(\mathcal F_{n-1}\)-measurable for \(n\geq1\).
This decomposition is unique up to almost-sure equality at every time
among decompositions with a martingale part, a predictable part and
initial predictable value zero.

The predictable part is nondecreasing almost surely exactly when
\(Y\) is a submartingale. In that case the increments \(d_n\)
may be chosen nonnegative everywhere, giving an everywhere
nondecreasing predictable version
\cite[thm.~4.3.2]{DurrettPTEAuthor2019}.
\index{Doob decomposition}
\end{hbtheorem}

\begin{hbtheorem}{thm:p15-bounded-increment-dichotomy}{proved=annex:p15.bounded-increment-dichotomy}{Convergence or two-sided unbounded oscillation}
If \(X\) is a martingale and
\(|X_{n+1}-X_n|\leq K<\infty\) almost surely for every \(n\),
where \(K\) is deterministic, then with probability one either
\(X_n\) has a finite limit or
\[
                      \liminf_nX_n=-\infty,\qquad
                      \limsup_nX_n=+\infty.
\]
No boundedness assumption on the integrable initial value is needed.
A bound that is merely finite separately on almost every sample path
does not suffice \cite[thm.~4.3.1]{DurrettPTEAuthor2019}.
\end{hbtheorem}

\begin{hbdefinition}{def:p15-adaptive-counts}{Events, their conditional probabilities and the compensated count}
Let \(B_n\in\mathcal F_n\), \(n\geq1\), and choose
\[
 p_n=E(\mathbf1_{B_n}\mid\mathcal F_{n-1})\in[0,1],\qquad
 C_n=\sum_{j=1}^n\mathbf1_{B_j},\qquad
 P_n=\sum_{j=1}^np_j,\qquad M_n=C_n-P_n,
\]
with all three sums zero at time zero.
The chosen probability versions are bounded and predictable.
Write \(P_\infty=\sum_{j\geq1}p_j\).
The symbol \(P_n\) here is a random compensating sum, whereas
unsubscripted \(P\) denotes the underlying probability measure.
\index{conditional Borel--Cantelli lemma}\index{compensated event count}
\end{hbdefinition}

\begin{hbtheorem}{thm:p15-conditional-borel-cantelli-source}{proved=annex:p15.conditional-borel-cantelli-source}{The source's convergence-or-oscillation count proof}
The compensated process \(M_n\) is a martingale with
\(|\Delta M_n|\leq1\), and
\[
                   \{B_n\text{ infinitely often}\}
                                  =\{P_\infty=\infty\}
                                          \quad\text{almost surely}.
\]
The full proof uses the preceding bounded-increment dichotomy.
It applies to arbitrary \(\mathcal F_0\), including the trivial
initial field stipulated in the source
\cite[ex.~4.3.3, thm.~4.3.4]{DurrettPTEAuthor2019}.
\end{hbtheorem}

\begin{hbtheorem}{thm:p15-summable-adapted-errors}{proved=annex:p15.summable-adapted-errors}{Summable allowances and nonnegative decrements}
Let \((Z_n,U_n,V_n)_{n\geq0}\) be nonnegative, integrable and
\(\mathcal F_n\)-measurable, with
\[
                 E(Z_{n+1}\mid\mathcal F_n)
                                 \leq Z_n-U_n+V_n.
\]
On the event \(\{\sum_nV_n<\infty\}\), the sequence \(Z_n\)
converges to a finite limit and \(\sum_nU_n<\infty\), almost surely.
Each allowance and decrement is known before the next step.
The case \(U_n=0\) supplies the full source's summable-error exercise;
no assumption \(E\sum_nV_n<\infty\) is imposed
\cite[ex.~4.3.3]{DurrettPTEAuthor2019}.
\end{hbtheorem}

\section{Predictable and realized squared variation}
\label{sec:p15-quadratic-variation-definitions}

\begin{hbdefinition}{def:p15-predictable-variation}{The variance process}
In the rest of the chapter, unless explicitly stated otherwise,
\(X\) is a real square-integrable martingale with \(X_0=0\).
For \(n\geq1\), choose finite nonnegative
\(\mathcal F_{n-1}\)-measurable versions
\[
 a_n=E((\Delta X_n)^2\mid\mathcal F_{n-1}),\qquad
 A_0=0,\quad A_n=\sum_{j=1}^na_j,\quad A_\infty=\sup_nA_n.
\]
The process \(A\), also denoted \(\langle X\rangle\), is the
\emph{predictable quadratic variation}. Its next value \(A_n\)
is already known at time \(n-1\).
The \emph{realized squared variation} is
\[
                  [X]_0=0,\qquad [X]_n=\sum_{j=1}^n(\Delta X_j)^2.
\]
It is adapted and nondecreasing, but need not be predictable.
Write \(X^*=\sup_{n\geq0}|X_n|\).
All assertions said to hold almost surely on an event \(E\)
mean that their failures intersect \(E\) in a null set.
\index{quadratic variation!predictable}\index{quadratic variation!realized}
\end{hbdefinition}

\begin{hbtheorem}{thm:p15-predictable-variation-identities}{proved=annex:p15.predictable-variation-identities}{Compensation and the discrete square identity}
The process \(A\) is predictable, nondecreasing and integrable
at every finite time. Both \(X_n^2-A_n\) and \([X]_n-A_n\)
are integrable martingales, and
\[
                 EA_n=EX_n^2=E([X]_n).
\]
The pathwise square identity is
\[
                 X_n^2=2\sum_{j=1}^nX_{j-1}\Delta X_j+[X]_n.
\]
Its stochastic sum is an integrable martingale even though its
predictable coefficients need not be bounded. The required product
integrability follows from square integrability.
\end{hbtheorem}

\begin{hbtheorem}{thm:p15-variation-stopping-transform}{proved=annex:p15.variation-stopping-transform}{Stopping and predictable transformation of the variance process}
For a stopping time \(T\), the stopped martingale \(X_{T\wedge n}\)
has predictable variation \(A_{T\wedge n}\).
For predictable coefficients \(H_j\), bounded by deterministic
constants at each finite time, the martingale
\[
                         Y_n=\sum_{j=1}^nH_j\Delta X_j
\]
has predictable variation
\[
                         \langle Y\rangle_n=\sum_{j=1}^nH_j^2a_j.
\]
The corresponding stopped realized variation is \([X]_{T\wedge n}\).
Writing \(A_T=\lim_nA_{T\wedge n}\) and
\([X]_T=\lim_n[X]_{T\wedge n}\), one also has
\[
                               E([X]_T)=EA_T
\]
as an extended nonnegative identity.
Even when \(T\) is finite almost surely, this alone does not
imply \(EX_T^2=EA_T\) at an unbounded time.
\end{hbtheorem}

\begin{hbtheorem}{thm:p15-quadratic-maximal-bound}{proved=annex:p15.quadratic-maximal-bound}{The square maximal bound}
One has
\[
                                  E(X^*)^2\leq4EA_\infty.
\]
If \(EA_\infty<\infty\), then \(X_n\) converges almost surely
and in \(L^2\) to \(X_\infty\), and
\(EX_\infty^2=EA_\infty\)
\cite[thm.~4.5.1]{DurrettPTEAuthor2019}.
\end{hbtheorem}

\begin{hbtheorem}{thm:p15-finite-variation-convergence}{proved=annex:p15.finite-variation-convergence}{Pathwise finite variance gives a finite limit}
The martingale \(X_n\) converges to a finite limit almost surely
on \(\{A_\infty<\infty\}\).
For each \(c>0\), the time
\[
                     T_c=\inf\{n\geq0:A_{n+1}>c\}
\]
is a stopping time, and its stopped variation satisfies
\[
                       A_{T_c\wedge n}\leq c\quad\text{for every }n.
\]
Stopping before the next variance increment is essential to this
bound \cite[thm.~4.5.2]{DurrettPTEAuthor2019}.
\end{hbtheorem}

\begin{hbtheorem}{thm:p15-bounded-jump-variation-dichotomy}{proved=annex:p15.bounded-jump-variation-dichotomy}{The converse when jumps have a fixed bound}
If additionally \(|\Delta X_n|\leq K<\infty\) for a deterministic
\(K\), then, up to null sets,
\[
 \{X_n\text{ has a finite limit}\}=\{A_\infty<\infty\}.
\]
On \(\{A_\infty=\infty\}\), its lower and upper limits are
\(-\infty\) and \(+\infty\), respectively.
The converse fails without the fixed jump bound.
\end{hbtheorem}

\section{Normalization and adaptive frequencies}
\label{sec:p15-variance-normalization}

\begin{hbtheorem}{thm:p15-variance-clock-normalization}{proved=annex:p15.variance-clock-normalization}{Normalization by a predictable variance scale}
Let \(f:[0,\infty)\to[1,\infty)\) be a finite nondecreasing
Borel function with
\[
                             I_f=\int_0^\infty f(t)^{-2}\,dt<\infty.
\]
Then
\[
                  \frac{X_n}{f(A_n)}\longrightarrow0
                                \quad\text{almost surely on }\{A_\infty=\infty\}.
\]
The transformed martingale
\(Y_n=\sum_{j=1}^n\Delta X_j/f(A_j)\) has
\(\langle Y\rangle_\infty\leq I_f\).
The proof retains the source's predictable denominator \(f(A_j)\)
and its pathwise Kronecker argument
\cite[thm.~4.5.3]{DurrettPTEAuthor2019}.
\end{hbtheorem}

\begin{hbtheorem}{thm:p15-variance-logarithmic-rate}{proved=annex:p15.variance-logarithmic-rate}{The logarithmic variance normalization}
For \(\epsilon>0\), use natural logarithms and define
\[
 f_\epsilon(t)=
 \begin{cases}
 \sqrt e,&0\leq t\leq e,\\
 \sqrt{\,t(\log t)^{1+\epsilon}\,},&t>e.
 \end{cases}
\]
This is finite, nondecreasing and at least one, and
\[
                           \int_0^\infty f_\epsilon(t)^{-2}\,dt
                                                    =1+\epsilon^{-1}.
\]
Consequently, on \(\{A_\infty=\infty\}\),
\[
                  \frac{X_n}{\sqrt{A_n(\log A_n)^{1+\epsilon}}}
                                   \longrightarrow0\quad\text{almost surely},
\]
where the displayed denominator is used once \(A_n>e\).
For independent centered increments of variances \(\sigma_j^2\),
\(A_n=\sum_{j=1}^n\sigma_j^2\).
This gives the source's nonidentically distributed extension
\cite[ex.~4.5.4]{DurrettPTEAuthor2019}.
\end{hbtheorem}

\begin{hbtheorem}{thm:p15-conditional-borel-cantelli-ratio}{proved=annex:p15.conditional-borel-cantelli-ratio}{The normalized conditional Borel--Cantelli law}
For the adaptive counts of \cref{def:p15-adaptive-counts},
\[
                         \langle M\rangle_n
                                     =\sum_{j=1}^np_j(1-p_j)\leq P_n.
\]
On \(\{P_\infty=\infty\}\),
\[
                                \frac{C_n}{P_n}\longrightarrow1
                                               \quad\text{almost surely}.
\]
The ratio is used only once \(P_n>0\), which holds eventually
on that event. On \(\{P_\infty<\infty\}\), only finitely many
\(B_n\) occur almost surely.
The proof separately treats finite variance with divergent \(P_n\)
and divergent variance; the first case cannot be discarded
\cite[thm.~4.5.5 and remark]{DurrettPTEAuthor2019}.
\end{hbtheorem}

\begin{hbdefinition}{def:p15-friedman-urn}{A two-colour urn with opposite-colour reinforcement}
Let \(g,r>0\), \(a\geq0\) and \(b>0\) be finite real weights.
Initially \(G_0=g\), \(R_0=r\). At draw \(n\geq1\), let \(I_n\)
be the indicator of green, chosen conditionally with probability
\(g_{n-1}=G_{n-1}/(G_{n-1}+R_{n-1})\).
Keep the existing weights, add \(a\) to the selected colour and
\(b\) to the opposite colour. Thus, with \(D_n=\sum_{j=1}^nI_j\),
\[
 \begin{split}
 G_n&=g+aD_n+b(n-D_n),\\
 R_n&=r+bD_n+a(n-D_n),\\
 t_n&=G_n+R_n=g+r+(a+b)n,\qquad
 g_n=G_n/t_n,\quad r_n=R_n/t_n.
 \end{split}
\]
Independent uniform variables realize the conditional draws.
Integer weights give the literal ball model; real weights give the
same specified probability scheme.
\index{Friedman urn}
\end{hbdefinition}

\begin{hbtheorem}{thm:p15-friedman-urn-limit}{proved=annex:p15.friedman-urn-limit}{The complete urn limit, including zero same-colour reinforcement}
For every parameter choice in the preceding definition,
\[
                                   g_n\longrightarrow\frac12
                                          \quad\text{almost surely}.
\]
Both colours are drawn infinitely often, and
\[
                        D_n/n\longrightarrow\frac12
                                          \quad\text{almost surely}.
\]
The source's cases \(a=b\), \(a>b\) and \(0<a<b\) retain their
full count-and-limsup proofs. The boundary \(a=0\) is supplied by
a quadratic-drift argument with summable allowances
\cite[ex.~4.5.6]{DurrettPTEAuthor2019}.
\end{hbtheorem}

\section{The square-root estimate and unbounded stopping}
\label{sec:p15-square-root-variation}

\begin{hbtheorem}{thm:p15-square-root-variation-bound}{proved=annex:p15.square-root-variation-bound}{The maximal first moment from square-root variance}
For the square-integrable martingale starting at zero,
\[
                                  EX^*\leq3E\sqrt{A_\infty}.
\]
This is an extended nonnegative inequality; in particular a finite
right side gives an integrable all-time maximum.
The proof supplies the entire source stopping, tail-integration and
nonnegative interchange argument
\cite[thm.~4.5.7]{DurrettPTEAuthor2019}.
\end{hbtheorem}

\begin{hbtheorem}{thm:p15-square-root-stopping}{proved=annex:p15.square-root-stopping}{A sufficient unbounded sampling criterion}
Let \(T\) be an almost-surely finite stopping time.
If \(E\sqrt{A_T}<\infty\), then
\[
                       E\sup_n|X_{T\wedge n}|\leq3E\sqrt{A_T},
 \qquad X_{T\wedge n}\longrightarrow X_T\text{ in }L^1,\qquad EX_T=0.
\]
For any stopping time \(S\leq T\) pointwise,
\[
                              E(X_T\mid\mathcal F_S)=X_S.
\]
If \(EA_T<\infty\), the stopped convergence is also in \(L^2\)
and \(EX_T^2=EA_T\).
For a nonzero square-integrable initial value, apply these statements
to \(X-X_0\); the expectation conclusion becomes \(EX_T=EX_0\).

For the fair-sign random walk \(S_0=0\), \(E\sqrt T<\infty\)
therefore implies \(ES_T=0\). Its first hitting time \(H\) of
\(-1\) is finite almost surely but satisfies \(E\sqrt H=\infty\)
\cite[ex.~4.5.8]{DurrettPTEAuthor2019}.
\end{hbtheorem}

\section{Calculated examples and exact limitations}
\label{sec:p15-quadratic-variation-boundaries}

\begin{hbtheorem}{thm:p15-variance-clock-examples}{proved=annex:p15.variance-clock-examples}{Finite limits, realized variation and different moment requirements}
\begin{enumerate}
\item Independent increments taking values \(n,-n,0\) with probabilities
      \(1/(2n^2),1/(2n^2),1-1/n^2\), respectively, give \(A_n=n\).
      Almost surely only finitely many increments are nonzero, so
      \(X_n\) is eventually constant and \([X]_\infty<\infty\),
      although \(A_\infty=\infty\).
\item Let \(N\) be independent of a fair-sign walk with \(S_0=0\) and known at time
      zero, with
      \[
                  P(N=k)=c\,k^{-7/4},\qquad
                  c=\left(\sum_{k=1}^\infty k^{-7/4}\right)^{-1}.
      \]
      For \(X_n=S_{N\wedge n}\), \(A_n=N\wedge n\),
      \(E\sqrt{A_\infty}<\infty\), but \(EA_\infty=\infty\).
      The terminal value has mean zero and finite first absolute
      moment, but infinite second moment.
\item A single fair-sign step followed by constant values has
      \(A_n=1\) for \(n\geq1\). For a function \(f\) in the
      normalization theorem, its limit divided by \(f(1)\)
      is a nonzero signed value. The infinite-variance event in
      \cref{thm:p15-variance-clock-normalization} is essential.
\end{enumerate}
\end{hbtheorem}

\begin{hbtheorem}{thm:p15-randomly-bounded-paths}{proved=annex:p15.randomly-bounded-paths}{A martingale with bounded sample paths and no limit}
Fix \(0<p<1/2\). Let \(\zeta_n\) be i.i.d. with probabilities
\(p,1-2p,p\) at \(-1,0,1\), and independently let
\(\eta_n\) be independent Bernoulli variables of probabilities
\(\epsilon_n=2^{-n}\). Define
\[
                   X_0=0,\qquad
                   X_n=(1-\eta_n)\zeta_n
                                +\eta_nX_{n-1}/\epsilon_n.
\]
In the filtration generated by these inputs, \(X\) is a
square-integrable martingale at every finite time and
\[
          E|X_n|=2p(n-1+2^{-n}),\qquad
          P(X_n=a)\longrightarrow P(\zeta_1=a)
                                \quad(a=-1,0,1).
\]
Almost surely \(\sup_n|X_n|<\infty\) and each of the values
\(-1,0,1\) occurs infinitely often. It has no almost-sure or
probability limit, and no deterministic common increment bound.
This supplies both source Exercises~4.3.1--2.
\end{hbtheorem}

\begin{hbtheorem}{thm:p15-urn-boundary-calculations}{proved=annex:p15.urn-boundary-calculations}{Why opposite-colour reinforcement matters}
If \(a=b>0\), the Friedman urn has the deterministic fractions
\[
                                g_n=\frac{g+an}{g+r+2an}.
\]
If instead \(b=0,a=1,g=r=1\), the same-colour P\'olya urn has
a bounded martingale fraction whose limit is uniform on \((0,1)\),
rather than equal to \(1/2\) almost surely.
After \(n\) draws its green count \(G_n\) is uniform on
\(\{1,\ldots,n+1\}\).

For adaptive counts, the deterministic events \(B_n=\Omega\)
give \(P_n=n\) but \(\langle M\rangle_n=0\).
More generally, let \(\Theta\in[0,1]\) be known at time zero and
independent of a sequence of mutually independent uniform inputs
\(U_n\in(0,1)\), and take
\[
                              B_n=\{U_n\leq\Theta/(n+1)\}.
\]
Then \(p_n=\Theta/(n+1)\), and, for \(n\geq2\),
\(C_n/\log n\to\Theta\) almost surely.
This includes the zero-\(\Theta\) event without dividing by it.
\end{hbtheorem}

\begin{hbtheorem}{thm:p15-decomposition-rate-boundaries}{proved=annex:p15.decomposition-rate-boundaries}{Predictability, initial values and an insufficient rate}
\begin{enumerate}
\item Let \(B\) be Bernoulli \((1/2)\), \(Y_0=0\), and \(Y_n=B\)
      for \(n\geq1\). It has the predictable drift
      \(D_0=0,D_n=1/2\) for \(n\geq1\), and martingale remainder
      \(M_0=0,M_n=B-1/2\) for \(n\geq1\).
      The alternative \(Y=0+Y\) has an adapted increasing second
      part which is not predictable. Dropping predictability
      destroys uniqueness.
\item A nonzero constant martingale has zero variance process.
      Thus the square and square-root maximal estimates require
      either \(X_0=0\) or the stated recentering.
\item For a fair-sign walk starting at zero, \(A_n=n\) and
      \(f(t)=\max(1,\sqrt t)\) is finite and nondecreasing, but
      \(\int_0^\infty f(t)^{-2}\,dt=\infty\).
      In fact
      \[
                            P(|S_n|>\sqrt{n/2})\geq1/12
                                                \quad(n\geq1),
      \]
      so the proposed normalized limit zero fails even in probability.
\end{enumerate}
\end{hbtheorem}

The density-martingale decomposition and full independent-product
criterion are supplied in \cref{thm:p15-density-lebesgue-decomposition,thm:p15-product-positive-tail,thm:p15-product-affinity-criterion}.
The general P\'olya beta limit and the source's two special-case
calculations are supplied in \cref{thm:p15-polya-beta-limit,thm:p15-polya-source-special-cases}.
The complete Section~4.8 optional-sampling and random-walk constructions
are supplied in \cref{sec:p15-optional-stopping-and-random-walks}.
